% GENERATED FILE. Edit canonical lessons, then run npm run generate:books.
% canonical-source-hash: fnv1a64:03cad2e7bb233349
% canonical-lessons: MR-C330-senapati, MR-C330-paydal, MR-C330-chavni, MR-C330-pahara-dene, MR-C330-jhunjne

\chapter{The Army in Camp}
\label{ch:mr-the-army-in-camp}
\hlchaptermodality{330}

\begin{chapteropening}
\textbf{By the end of this chapter:} \emph{I can talk about commanders, foot soldiers and guards.}

Complete the last lesson of chapter 330: I can talk about commanders, foot soldiers and guards.
\end{chapteropening}

\section[\texorpdfstring{\mr{सेनापती} (senāpatī)}{सेनापती (senāpatī)}]{\mr{सेनापती} (senāpatī) — a commander}

\label{lesson:MR-C330-senapati}



\begin{quote}
Before the new one: say the Marathi for a cannon, then the Marathi for cavalry.
\end{quote}

\subsection*{You'll want to know: \mr{सेनापती}}
\textbf{\mr{सेनापती}} — \emph{senāpatī} — \textquotedblleft{}a commander\textquotedblright{}.

\begin{quote}
\mr{सेनापतीने} \mr{सैन्याला} \mr{थांबायला} \mr{सांगितलं}.
\end{quote}

\begin{quote}
\emph{senāpatīne sainyālā thāmbāylā sāṅgitlaṁ.}
\end{quote}

\begin{quote}
The commander told the army to halt.
\end{quote}

Masculine.

\mr{सेना}, an army, + \mr{पती}, a lord.

\begin{cousinweb}
\textbf{\mr{सेनापती}} (\emph{senāpatī}) is a tatsama, taken from Sanskrit \emph{senāpati}, 'the lord of an army, a general' (Monier-Williams).
\end{cousinweb}

\begin{grammarlens}[title={What you've built}]
One more word for telling what happened.
\end{grammarlens}

\subsection*{Your turn}
\begin{itemize}
\raggedright
  \item \emph{Say it:} \emph{senāpatī}
  \item \emph{Say it:} \emph{senāpatī}, once more
  \item \emph{Say it:} \emph{senāpatīne sainyālā thāmbāylā sāṅgitlaṁ.} — The commander told the army to halt.
  \item \emph{Recall:} say the Marathi for a cannon, then the Marathi for cavalry, then say \emph{senāpatī} again
\end{itemize}

\begin{tcolorbox}[breakable,colback=teal!4,colframe=teal!35!black,title={Before you move on}]
What does \mr{सेनापती} mean? (\textquotedblleft{}A commander\textquotedblright{}.) Say it once more.
\end{tcolorbox}
\section[\texorpdfstring{\mr{पायदळ} (pāydaḷ)}{पायदळ (pāydaḷ)}]{\mr{पायदळ} (pāydaḷ) — infantry}

\label{lesson:MR-C330-paydal}



\begin{quote}
Before the new one: say the Marathi for cavalry, then the Marathi for a commander.
\end{quote}

\subsection*{You'll want to know: \mr{पायदळ}}
\textbf{\mr{पायदळ}} — \emph{pāydaḷ} — \textquotedblleft{}infantry\textquotedblright{}.

\begin{quote}
\mr{पायदळ} \mr{घोडदळाच्या} \mr{मागे} \mr{चालत} \mr{होतं}.
\end{quote}

\begin{quote}
\emph{pāydaḷ ghoḍdaḷācyā māge cālat hotaṁ.}
\end{quote}

\begin{quote}
The infantry was marching behind the cavalry.
\end{quote}

Neuter.

\mr{पाय}, a foot, + \mr{दळ}, a troop.

\begin{cousinweb}
\textbf{\mr{पायदळ}} (\emph{pāydaḷ}) is built from \emph{pāy}, 'a foot', and \emph{daḷ}, 'a troop', from Sanskrit \emph{dala} (Monier-Williams): the troop on foot.
\end{cousinweb}

\begin{grammarlens}[title={What you've built}]
One more word for telling what happened.
\end{grammarlens}

\subsection*{Your turn}
\begin{itemize}
\raggedright
  \item \emph{Say it:} \emph{pāydaḷ}
  \item \emph{Say it:} \emph{pāydaḷ}, once more
  \item \emph{Say it:} \emph{pāydaḷ ghoḍdaḷācyā māge cālat hotaṁ.} — The infantry was marching behind the cavalry.
  \item \emph{Recall:} say the Marathi for cavalry, then the Marathi for a commander, then say \emph{pāydaḷ} again
\end{itemize}

\begin{tcolorbox}[breakable,colback=teal!4,colframe=teal!35!black,title={Before you move on}]
What does \mr{पायदळ} mean? (\textquotedblleft{}Infantry\textquotedblright{}.) Say it once more.
\end{tcolorbox}
\section[\texorpdfstring{\mr{छावणी} (chāvṇī)}{छावणी (chāvṇī)}]{\mr{छावणी} (chāvṇī) — an army camp; a camp}

\label{lesson:MR-C330-chavni}



\begin{quote}
Before the new one: say the Marathi for a commander, then the Marathi for infantry.
\end{quote}

\subsection*{You'll want to know: \mr{छावणी}}
\textbf{\mr{छावणी}} — \emph{chāvṇī} — \textquotedblleft{}an army camp; a camp\textquotedblright{}.

\begin{quote}
\mr{नदीकाठी} \mr{सैन्याची} \mr{छावणी} \mr{पडली} \mr{होती}.
\end{quote}

\begin{quote}
\emph{nadīkāṭhī sainyācī chāvṇī paḍlī hotī.}
\end{quote}

\begin{quote}
The army's camp had been pitched on the riverbank.
\end{quote}

Feminine.

\mr{छावणी} \mr{पडणे}, a camp is pitched.

\begin{cousinweb}
\textbf{\mr{छावणी}} (\emph{chāvṇī}) is the Marathi word for an army camp; its origin is not settled.
\end{cousinweb}

\begin{grammarlens}[title={What you've built}]
One more word for telling what happened.
\end{grammarlens}

\subsection*{Your turn}
\begin{itemize}
\raggedright
  \item \emph{Say it:} \emph{chāvṇī}
  \item \emph{Say it:} \emph{chāvṇī}, once more
  \item \emph{Say it:} \emph{nadīkāṭhī sainyācī chāvṇī paḍlī hotī.} — The army's camp had been pitched on the riverbank.
  \item \emph{Recall:} say the Marathi for a commander, then the Marathi for infantry, then say \emph{chāvṇī} again
\end{itemize}

\begin{tcolorbox}[breakable,colback=teal!4,colframe=teal!35!black,title={Before you move on}]
What does \mr{छावणी} mean? (\textquotedblleft{}An army camp; a camp\textquotedblright{}.) Say it once more.
\end{tcolorbox}
\section[\texorpdfstring{\mr{पहारा} \mr{देणे} (pahārā deṇe)}{पहारा देणे (pahārā deṇe)}]{\mr{पहारा} \mr{देणे} (pahārā deṇe) — to stand guard}

\label{lesson:MR-C330-pahara-dene}



\begin{quote}
Before the new one: say the Marathi for infantry, then the Marathi for an army camp.
\end{quote}

\subsection*{You'll want to know: \mr{पहारा} \mr{देणे}}
\textbf{\mr{पहारा} \mr{देणे}} — \emph{pahārā deṇe} — \textquotedblleft{}to stand guard\textquotedblright{}.

\begin{quote}
\mr{दोन} \mr{सैनिकांनी} \mr{रात्रभर} \mr{दारावर} \mr{पहारा} \mr{दिला}.
\end{quote}

\begin{quote}
\emph{don sainikānnī rātrabhar dārāvar pahārā dilā.}
\end{quote}

\begin{quote}
Two soldiers stood guard at the gate all night.
\end{quote}

Past, with -\mr{ने} (he …): \mr{त्याने} \mr{पहारा} \mr{दिला}, \emph{tyāne pahārā dilā}.

The place guarded takes -\mr{वर}; \mr{पहारा} is masculine, so the past is \mr{पहारा} \mr{दिला}.

\begin{cousinweb}
\textbf{\mr{पहारा} \mr{देणे}} (\emph{pahārā deṇe}) is built from \emph{pahārā}, 'a watch, a guard', of the family of Sanskrit \emph{prahara}, 'a watch of about three hours' (CDIAL 8900), and \emph{deṇe}, 'to give'.
\end{cousinweb}

\begin{grammarlens}[title={What you've built}]
One more word for telling what happened.
\end{grammarlens}

\subsection*{Your turn}
\begin{itemize}
\raggedright
  \item \emph{Say it:} \emph{pahārā deṇe}
  \item \emph{Say it:} \emph{pahārā deṇe}, once more
  \item \emph{Say it:} \emph{don sainikānnī rātrabhar dārāvar pahārā dilā.} — Two soldiers stood guard at the gate all night.
  \item \emph{Recall:} say the Marathi for infantry, then the Marathi for an army camp, then say \emph{pahārā deṇe} again
\end{itemize}

\begin{tcolorbox}[breakable,colback=teal!4,colframe=teal!35!black,title={Before you move on}]
What does \mr{पहारा} \mr{देणे} mean? (\textquotedblleft{}To stand guard\textquotedblright{}.) Say it once more.
\end{tcolorbox}
\section[\texorpdfstring{\mr{झुंजणे} (jhuñjṇe)}{झुंजणे (jhuñjṇe)}]{\mr{झुंजणे} (jhuñjṇe) — to battle; to fight hard}

\label{lesson:MR-C330-jhunjne}



\begin{quote}
Before the new one: say the Marathi for an army camp, then the Marathi for to stand guard.
\end{quote}

\subsection*{You'll want to know: \mr{झुंजणे}}
\textbf{\mr{झुंजणे}} — \emph{jhuñjṇe} — \textquotedblleft{}to battle; to fight hard\textquotedblright{}.

\begin{quote}
\mr{मूठभर} \mr{सैनिक} \mr{मोठ्या} \mr{सैन्याशी} \mr{झुंजले}.
\end{quote}

\begin{quote}
\emph{mūṭhbhar sainik moṭhyā sainyāśī jhuñjle.}
\end{quote}

\begin{quote}
A handful of soldiers battled a large army.
\end{quote}

Past (he …): \mr{तो} \mr{झुंजला}, \emph{to jhuñjlā}.

The opponent takes -\mr{शी}: \mr{शत्रूशी} \mr{झुंजणे}.

\begin{cousinweb}
\textbf{\mr{झुंजणे}} (\emph{jhuñjṇe}) is a tadbhava: CDIAL lists Marathi \emph{jhuñjṇe}, 'to fight', under Sanskrit \emph{yudhyate}, 'fights against' (CDIAL 10502).
\end{cousinweb}

\begin{grammarlens}[title={What you've built}]
One more word for telling what happened.
\end{grammarlens}

\subsection*{Your turn}
\begin{itemize}
\raggedright
  \item \emph{Say it:} \emph{jhuñjṇe}
  \item \emph{Say it:} \emph{jhuñjṇe}, once more
  \item \emph{Say it:} \emph{mūṭhbhar sainik moṭhyā sainyāśī jhuñjle.} — A handful of soldiers battled a large army.
  \item \emph{Recall:} say the Marathi for an army camp, then the Marathi for to stand guard, then say \emph{jhuñjṇe} again
\end{itemize}

\begin{tcolorbox}[breakable,colback=teal!4,colframe=teal!35!black,title={Before you move on}]
What does \mr{झुंजणे} mean? (\textquotedblleft{}To battle; to fight hard\textquotedblright{}.) Say it once more.
\end{tcolorbox}
